% ============================================================================
% CHAPTER 9: Tools Anyone Can Use
% A practical survey of annotation platforms, HITL services, and how to
% get started with real data
% ============================================================================
% Source: editions/winter/chapters/ch9/Ch9_final.md
% ============================================================================

\label{chap:ch09}

\begin{storyhook}
  \textit{A practical survey of annotation platforms, \hitl services, and how to get started with real data.}
\end{storyhook}

% ============================================================================
\section{The Label Studio Moment}
\label{sec:ch09-label-studio-moment}
% ============================================================================

Somewhere in 2021, a solo developer named Nikolai Litvak published a blog post titled ``How I labeled 50,000 images in a weekend.'' The post described his experience using Label Studio, an open-source annotation tool he had built and released a year earlier. He had expected it to be used by machine learning researchers at well-funded labs. Instead, he found his community forums filling with posts from a different audience: non-profits trying to label satellite images of deforestation, a hospital system annotating medical records, a small news organization classifying articles for a study of media coverage.

None of these organizations had data science teams in the traditional sense. None of them could afford enterprise annotation platforms. But they had data, they had a task, and they had people willing to do the labeling work. What they needed was infrastructure.

This is the story of the past decade in \hitl tooling. The infrastructure---the platforms, the pipelines, the APIs, the hosted services---has become accessible to organizations that would have been entirely locked out five years earlier. The engineering sophistication required to build a human-in-the-loop system from scratch has dropped dramatically. And with it, the range of people and institutions who can participate in building \ai systems that use human judgment has expanded.

The tools are not magic. They have real limitations, and choosing the wrong tool for a task can cost more in rework than building the right thing from scratch. But the landscape is rich enough now that the question ``what tool should I use?'' has real answers.

% ============================================================================
\section{What Annotation Tools Actually Do}
\label{sec:ch09-what-tools-do}
% ============================================================================

Before surveying the landscape, a quick reminder of what these tools are for.

An annotation tool provides:
\begin{itemize}
  \item A way to present data items (images, text, audio, video, structured records) to reviewers
  \item A way for reviewers to record their judgments in structured form
  \item A way to store those judgments and export them for use in training
  \item Usually: some way to manage the workflow---who reviews what, when, quality control
\end{itemize}

That is the core. Advanced tools add:
\begin{itemize}
  \item Agreement tracking and quality metrics
  \item Active learning queues (routing cases by model uncertainty)
  \item Model-assisted annotation (the model pre-labels, humans correct)
  \item Analytics dashboards
  \item Multi-reviewer reconciliation
  \item Team management and access control
\end{itemize}

The choice of tool depends on your data type, your task complexity, your team size, your budget, and how much customization you need.

% ============================================================================
\section{Open-Source Options}
\label{sec:ch09-open-source}
% ============================================================================

\paragraph{Label Studio} is the most widely used open-source annotation platform as of 2024. It supports all major data types: text, images, audio, video, time series, and HTML. The configuration is JSON-based and relatively accessible. It has a Python SDK and REST API, active community development, and a commercial cloud offering for teams that want hosted infrastructure.

Best for: research projects, small-to-medium teams with technical capacity, organizations with privacy requirements that prevent cloud hosting.

Limitations: setup requires some engineering; advanced features (active learning integration, enterprise workflow management) require more work.

\paragraph{Argilla} (formerly Rubrix) focuses specifically on NLP annotation---text classification, named entity recognition, question answering, and evaluation of language model outputs. It integrates well with the Hugging Face ecosystem and has strong support for \rlhf-style feedback collection. It is notably good for \llm evaluation use cases.

Best for: \nlp teams, \llm evaluation, RAG pipeline improvement.

\paragraph{CVAT (Computer Vision Annotation Tool)} is purpose-built for computer vision tasks. Originally developed by Intel, it supports bounding boxes, polygons, keypoints, and semantic segmentation. It is the tool of choice for many computer vision annotation projects.

Best for: image and video annotation, object detection, segmentation tasks.

\paragraph{Prodigy} is a commercial-licensed but developer-friendly tool from the creators of spaCy. It is designed for fast, scriptable annotation workflows with a particular emphasis on active learning integration. Prodigy's model-in-the-loop approach, where the model pre-annotates and humans correct, can dramatically reduce annotation time on \nlp tasks.

Best for: \nlp tasks where speed is critical and you have some technical capacity.

% ============================================================================
\section{Cloud Platform Approaches}
\label{sec:ch09-cloud-platforms}
% ============================================================================

For organizations that want managed infrastructure rather than self-hosted tooling, several cloud platforms offer \hitl services.

\paragraph{Scale AI} is the most prominent enterprise annotation platform. It offers managed annotation workforces, quality controls, domain-specific pipelines (medical, legal, autonomous vehicle), and an API for integrating annotation into \ml pipelines. Scale's managed workforce handles the recruitment, training, and quality management of annotators---a significant operational burden for organizations that do not want to manage their own annotation team.

Best for: large-scale annotation projects, organizations without in-house annotation capacity, high-stakes domains requiring managed quality control.

\paragraph{Amazon SageMaker Ground Truth} integrates annotation directly into AWS's \ml toolchain. It offers both Amazon Mechanical Turk (crowdsourced) and private workforce (your own team) options. It includes built-in active learning through Ground Truth Plus, which automatically sends uncertain cases to human reviewers and uses confirmed labels to improve the model.

Best for: organizations already in the AWS ecosystem, projects that benefit from tight \ml pipeline integration.

\paragraph{Google Vertex AI Data Labeling} is Google Cloud's managed annotation service. Similar in positioning to SageMaker Ground Truth but within GCP. Tight integration with Vertex AI for active learning loops.

\paragraph{Labelbox} is a mid-market platform with strong support for complex multi-step annotation projects, quality assurance workflows, and a notable focus on video annotation.

Best for: medium-to-large organizations with complex annotation workflows, teams that need strong quality assurance tooling.

% ============================================================================
\section{Crowdsourcing vs.\ Expert Labeling}
\label{sec:ch09-crowd-vs-expert}
% ============================================================================

This is one of the most consistently misunderstood trade-offs in annotation work.

\paragraph{Crowdsourced annotation} (Amazon Mechanical Turk, Scale AI's managed workforce, similar platforms) provides large-scale, fast, relatively low-cost annotation. The workforce is diverse and large. The limitation: for any task with significant domain requirements---medical text, legal documents, specialized imagery, culturally specific content---crowdworkers without that domain expertise produce noisy labels. Noise in training data is not random; it clusters around the hard cases that require domain judgment. These are often the most important cases.

\paragraph{Expert annotation} (domain specialists---doctors, lawyers, linguists, security analysts) produces high-quality labels on hard cases. The limitation: expensive, slow, difficult to scale, and hard to find. Expert annotators are also not immune to bias---their expertise brings systematic viewpoints that may not reflect the diversity of real-world cases \autocite{davani2022dealing}.

The practical resolution for most organizations: use the task to determine who labels it. Objectively determinable tasks (does this image contain a dog?) can use crowdsourcing successfully. Subjectively determinable tasks requiring domain judgment (does this symptom presentation suggest a particular diagnosis?) require experts. For tasks in between, consider a two-stage approach: crowd for initial filtering, experts for ambiguous cases.

Chapter~7 covers inter-annotator agreement in depth. Here, the practical tooling note: most annotation platforms support measuring inter-annotator agreement automatically (Krippendorff's $\alpha$, Cohen's $\kappa$), and using agreement scores to route low-agreement items to additional review.

% ============================================================================
\section{No-Code Options and Their Limits}
\label{sec:ch09-no-code}
% ============================================================================

Several platforms offer annotation workflows without any coding required.

\paragraph{Docanno} and \paragraph{Brat} are simple text annotation tools requiring minimal technical setup. They are effective for single annotators or small teams on text classification tasks, but do not scale well.

\paragraph{Roboflow} provides a no-code interface for computer vision annotation with integrated training and deployment. For image classification and object detection projects with small-to-medium scale, it is a complete pipeline with minimal technical friction.

\paragraph{Datasaur} offers a no-code text annotation platform with good \nlp support and collaboration features.

The honest assessment of no-code options: they work well for small, well-defined tasks. They hit ceilings quickly when you need active learning integration, custom label schemas, privacy requirements, or at-scale annotation. They are a useful starting point, not a long-term architecture.

% ============================================================================
\section{When to Build vs.\ Buy}
\label{sec:ch09-build-vs-buy}
% ============================================================================

The build vs.\ buy decision in \hitl tooling is more nuanced than in most software contexts because the annotation workflow is coupled to the \ml pipeline in ways that create vendor lock-in risks.

\paragraph{Buy when:}
\begin{itemize}
  \item Your task fits a standard annotation template (object detection, text classification, NER)
  \item You need scale quickly and do not have engineering capacity
  \item Data privacy is not a concern that rules out cloud hosting
  \item Your workflow requirements are stable (you know what you are building)
\end{itemize}

\paragraph{Build when:}
\begin{itemize}
  \item Your annotation task is non-standard or highly domain-specific
  \item You have strong privacy requirements that rule out external processing
  \item Your \hitl workflow needs deep integration with your \ml pipeline
  \item You have long-term product vision that requires full workflow control
\end{itemize}

\begin{technote}[The Hybrid Reality]
Most organizations end up using an off-the-shelf annotation tool (open-source or commercial) as the user interface layer, while building custom integration with their \ml infrastructure for routing, feedback, and monitoring. The annotation UX is not where the architectural moat is. The feedback integration layer is.
\end{technote}

% ============================================================================
\section{Getting Started Practically}
\label{sec:ch09-getting-started}
% ============================================================================

A realistic path to a working \hitl system for a solo developer or small team:

\paragraph{Week~1: Define the task precisely.} What is the input? What label is a reviewer assigning? What does ``correct'' look like? Write the annotation guidelines before you pick the tool. If you cannot write clear annotation guidelines, the task is not defined enough to build a system for.

\paragraph{Week~2: Collect a small sample and try labeling it yourself.} 100--200 examples. Where do you hesitate? Those are your hard cases. How long does each label take? That determines your reviewer throughput estimate.

\paragraph{Week~3: Set up Label Studio or Argilla.} Get a working annotation interface. Have a second person try to use it without help. Fix the parts they cannot figure out.

\paragraph{Week~4: Train a baseline model on your first 200 labels.} Even a simple model. Use it to get uncertainty estimates. Route the bottom third by confidence to your review queue.

\paragraph{Month~2: Iterate.} Review 200~more cases. Retrain. Check whether the low-confidence cases the model routed were actually harder. If they were---your uncertainty detection is working. If not---debug your calibration.

This is not glamorous. It is what building a real \hitl system looks like.

\begin{trythis}
Visit the Label Studio GitHub repository and browse the example templates. Pick one task that matches something you know about. Spend 15~minutes imagining what annotation guidelines you would write for that task. What is easy to write? What is hard? What would you need to know before you could write clear instructions? This exercise---writing annotation guidelines before touching any code---is the most important diagnostic for whether a \hitl project is ready to build.
\end{trythis}

% ============================================================================
\section*{Chapter Summary}
\addcontentsline{toc}{section}{Chapter Summary}
% ============================================================================

\begin{keyconcept}[Key Concepts]
\begin{itemize}
  \item The annotation tooling landscape ranges from open-source (Label Studio, Argilla, CVAT) to cloud platforms (Scale AI, SageMaker Ground Truth) to no-code options
  \item Crowdsourced vs.\ expert annotation: match the annotator expertise level to the task's domain requirements
  \item No-code tools hit ceilings quickly; the right long-term architecture often combines off-the-shelf annotation UX with custom feedback integration
  \item Build vs.\ buy depends on task novelty, privacy requirements, and long-term pipeline coupling
  \item The most important first step is writing annotation guidelines---before choosing a tool
\end{itemize}
\end{keyconcept}

\subsection*{Key Examples}

\begin{description}
  \item[Label Studio] The tool that democratized annotation for non-\ml organizations---used by non-profits, hospitals, and news organizations
  \item[\rlhf-focused tools like Argilla] Specialized for \llm evaluation and preference collection
  \item[Active learning in SageMaker Ground Truth] Architectural integration of uncertainty routing in a managed cloud service
\end{description}

\subsection*{Five Dimensions Check}

\begin{itemize}
  \item \emph{Feedback Integration} (Dimension~5): Tools are the infrastructure for closing the loop
  \item \emph{Intervention Design} (Dimension~2): The annotation interface is the intervention design implementation
\end{itemize}

\bigskip
\begin{center}
\rule{0.5\textwidth}{0.4pt}
\end{center}
\bigskip

\noindent\textit{Next chapter: language is the hardest thing to teach a computer, and it is also where the most interesting \hitl work is happening right now.}
